\chapter{Commentary on \PYS\ 1.23}
\renewcommand{\rightmark}{\MakeUppercase{1.23]}}
\thispagestyle{plain}
\label{trl:on1.23}

\bht{Does [the acquisition\pratikap{p23_1} of \samadhi] come extremely near just as a result of this?; is there another means for the attainment of that [\samadhi], or is there not?}\footnote{This introductory question in the \bhasya\ follows discussions in the preceding. However, as to what exactly the \bhasya\ is referring to as ``this'' is difficult to ascertain due to textual problems of the \PYS. Apparently the text of \PYS\ 1.21–22 that the author of our text knew was significantly different from what we can know from the \PYS's manuscript tradition. It is even probable that he did not consider what we consider as \sutra s 1.21--22 as \sutra s. See the relevant portion of my earlier edition \citep[230–1]{harimoto:diss} and Maas' critical edition of the \PYS, as well as his reconstruction of the text known to the author of the YVi based on my earlier edition \citep[32–33, 103, 141]{maas:2006}.  \citet{bronkhorst:patanjali+ys} theorizes that the author of the \bhasya\ was the compiler of the \sutra s based on this part of the \PYS. He did not consider the possibility that there might have been a significantly different version of the text. All the same, it can still be reasonably assumed that \PYS\ 1.21–22 classifies yogins and teaches that the yogin who works the hardest and/or the most talented achieves \samadhi\ the quickest. Hence the masculine singular word modified by \emph{āsannataraḥ} should be \emph{samādhilābhaḥ} (found in our text, too, to be modified by the word \emph{āsannaḥ}) and the word \emph{etasmāt} should refer to the whole passage that immediately precedes where one cause of quick achievement of \samadhi\ was taught.}  [The \Sutrakara] states that there is another means:

\begin{quotation}
\st{Or\pratikap{p23_2} by contemplation of \Isvara. (\rYS{1.23}{1.23})}
\end{quotation}

[The \Sutrakara]\pratikap{p23_3} will explain the meaning of the word ``\Isvara'' in the next [\sutra].\footnote{In the YVi the word \emph{uttara}, when referring to a portion of the root text, is used to refer to the next \sutra.  The next \sutra\ indeed defines \Isvara.  Accordingly, I consider the unspecified subject of the sentence to be the author of the \sutra.}  Here [the \Bhasyakara]  explains contemplation: [\Isvara,] \bht{Influenced by distinguished devotion}--- stimulated toward the inclination to helping---\bht{helps him}, the exuberantly devoted yogin, \bht{by [his] mere wish}, because the intentions of \Paramesvara\ [become] true\footnote{The characterization of the supreme deity/principle, \satyasamkalpa\ (the one whose intentions/plans become true), is, as \Sankara\ in \rBSBh{1.2.02}{1.2.2} says, an old one.} effortlessly. \bht{The attainment of \samadhi\ and [its] fruit become extremely near as a result of his help as well.}

